\begin{center}
\textsc{Gauss} an \textsc{Olbers}.\quad Göttingen, 9.~Oktober 1825.
\end{center}

\par\noindent\hspace*{\parindent}\mbox{.\ .\ .\ .\ .\ .}\quad Im Hinblick auf mein künftiges Werk über Höhere Geodäsie habe ich in diesen Tagen angefangen, einen (sehr) kleinen Teil dessen, was die gekrümmten Flächen betrifft, in Gedanken etwas zu ordnen.
Ich sehe aber, dass ich wegen der Eigenart meiner gesamten Darstellung um des Zusammenhangs willen \textls[120]{sehr weit ausholen} muss, sodass ich sogar meine Auffassung von den Krümmungsradien ebener Kurven vorausschicken muss.
Ich frage mich inzwischen fast, ob es ratsamer wäre, einen Teil dieser Lehren, der rein geometrisch (in analytischer Form) ist und Neues mit Bekanntem in neuer Form verbindet, zunächst eigens auszuarbeiten, ihn vielleicht aus dem Werk herauszulösen und als eine oder zwei Abhandlungen in unsere Commentationen aufzunehmen.
Vorerst kann ich jedoch die Form der Veröffentlichung offenlassen und einstweilen damit fortfahren, es zu Papier zu bringen.
\par\nobreak\noindent\hspace*{\parindent}\mbox{.\ .\ .\ .\ .\ .}\par
